% SSSUP-FONTE-STORICA: fonti-locali/originali/2014-2020.zip
% Anno 2017, esercizio 1. Testo storico; figura TikZ ricostruita in PSTricks.
% Verificare editorialmente la trasposizione prima di considerarla definitiva.
\renewcommand{\Anno}{2017}
\renewcommand{\NumeroProblema}{1}

\renewcommand{\Testo}{
In un'urna ci sono 100 palline, alcune bianche, le altre nere. Si sa che:
\begin{itemize}
\item il numero delle palline bianche è maggiore o uguale di quello delle nere;
\item se si estraggono contemporaneamente due palline dall'urna, la probabilità che siano entrambe dello stesso colore è uguale a quella che siano di colore diverso.
\end{itemize}

\begin{enumerate}
\item[(a)] Determinare il numero delle palline bianche.
\item[(b)] Se si sostituisce a 100 un intero $N$, discutere il problema precedente al variare di $N$.
\end{enumerate}
}

\renewcommand{\Risposta}{
% SSSUP-BOZZA-AI: Risposta; autore=ChatGPT; data=2026-11-10; stato=DA-VERIFICARE
% SSSUP-PROVENIENZA: ChatGPT (OpenAI), sessione ProblemiSSSUP del 2026-11-10; derivato dalla traccia canonica; revisione umana pendente.
(a) Ci sono $55$ palline bianche e $45$ nere.

(b) Per $N\ge2$ la condizione \`e realizzabile se e solo se $N=k^2$
per un intero $k\ge2$; in tal caso la composizione \`e univocamente
\[
B=\frac{k^2+k}{2},\qquad N-B=\frac{k^2-k}{2},
\]
dove $B$ indica il numero di palline bianche.
}
\renewcommand{\Aiutino}{
% SSSUP-BOZZA-AI: Aiutino; autore=ChatGPT; data=2026-11-10; stato=DA-VERIFICARE
% SSSUP-PROVENIENZA: ChatGPT (OpenAI), sessione ProblemiSSSUP del 2026-11-10; derivato dalla traccia canonica; revisione umana pendente.
Indicare con $B$ e $C$ i numeri di palline bianche e nere.
Calcolare in due modi il numero di coppie non ordinate di palline:
quelle dello stesso colore e quelle di colore diverso.
Dall'uguaglianza delle probabilit\`a si ottiene una condizione
aritmetica molto semplice sulla differenza $B-C$.
}
\renewcommand{\SoluzioneUfficiale}{}
\renewcommand{\SoluzioneRedazionale}{
% SSSUP-BOZZA-AI: SoluzioneRedazionale; autore=ChatGPT; data=2026-11-10; stato=DA-VERIFICARE
% SSSUP-PROVENIENZA: ChatGPT (OpenAI), sessione ProblemiSSSUP del 2026-11-10; derivato dalla traccia canonica; revisione umana pendente.
Scriviamo $B$ per il numero di palline bianche e $C=N-B$
per quello delle palline nere, con $B\ge C$.
Supponiamo $N\ge2$, affinch\'e l'estrazione di due palline sia possibile.
Le coppie non ordinate di palline distinte sono $\binom N2$ ed hanno
tutte la stessa probabilit\`a di essere estratte.
Le coppie dello stesso colore sono
$\binom B2+\binom C2$; quelle di colore diverso sono $BC$.
L'uguaglianza delle probabilit\`a equivale pertanto a
\[
\binom B2+\binom C2=BC,
\]
ossia, moltiplicando per $2$,
\[
 B(B-1)+C(C-1)=2BC.
\]
Riordinando e usando $B+C=N$ otteniamo
\[
 (B-C)^2=B+C=N.
\]

\textbf{(a)} Se $N=100$, siccome $B\ge C$, deve essere $B-C=10$.
Insieme a $B+C=100$ segue $B=55$ e $C=45$.

\textbf{(b)} Per un intero $N\ge2$, la condizione impone che
$N=k^2$ sia un quadrato perfetto, con $k=B-C\ge0$.
Viceversa, se $N=k^2$ con $k\ge2$, i numeri
\[
 B=\frac{k^2+k}{2},\qquad C=\frac{k^2-k}{2}
\]
sono interi non negativi, hanno somma $N$, differenza $k$ e soddisfano
l'uguaglianza delle probabilit\`a.
La soluzione \`e unica perch\'e sono assegnati sia la somma sia la
differenza. Per $N\ge2$ non quadrato non esiste alcuna composizione.
Il caso $N=1$ \`e escluso, perch\'e non si possono estrarre due palline.
}
\renewcommand{\SoluzioniAlternative}{}
\renewcommand{\Osservazioni}{}
\renewcommand{\Commenti}{}
